\section{Managerial and modeling implications}

The proposed framework demonstrates that integrating fractional-order dynamics, stochastic demand, carbon-trading mechanisms, and coherent risk management provides a more comprehensive representation of sustainable transportation planning under uncertainty than conventional memoryless formulations. The numerical analyses further highlight how these modeling components jointly influence transportation decisions, environmental performance, and system robustness.

The sensitivity analyses reveal that the fractional-order parameter significantly affects optimal transportation policies. Lower fractional orders promote more conservative decisions by encouraging preventive shipments and resource allocation, thereby reducing shortages and emergency interventions. These findings indicate that explicitly accounting for long-memory effects improves the representation of transportation systems characterized by persistent operational conditions, where the influence of past events cannot be neglected.

The computational results also demonstrate the strong interaction between environmental regulation and operational planning. Carbon allowances influence not only total emissions but also routing decisions, operating costs, and permit purchasing strategies. Restrictive carbon budgets increase compliance costs, whereas more flexible allocations improve operational adaptability. These observations confirm that carbon-trading mechanisms should be treated as endogenous decision components rather than external environmental constraints.

The risk analysis further shows that incorporating CVaR enhances the robustness of transportation plans. Increasing the degree of risk aversion consistently reduces exposure to extreme scenarios and service disruptions while producing only a moderate increase in expected transportation cost. This behavior illustrates the inherent trade-off between economic efficiency and operational resilience, confirming the suitability of CVaR as a coherent risk measure for sustainable transportation planning.

From a modeling perspective, the proposed formulation extends existing stochastic transportation models by combining fractional memory, environmental regulation, and coherent risk management within a unified optimization framework. The numerical experiments demonstrate that these three components are strongly interdependent and jointly determine transportation cost, carbon emissions, shortages, and service quality. Consequently, neglecting any one of these dimensions may lead to an incomplete characterization of transportation systems operating under uncertainty.

Overall, the proposed framework provides a unified decision-support methodology capable of simultaneously addressing long-memory effects, demand uncertainty, environmental constraints, and risk management. Beyond improving the realism of the mathematical formulation, the proposed approach establishes a flexible foundation for future developments involving large-scale transportation networks, dynamic carbon markets, real-time demand forecasting, and data-driven optimization strategies.